% 链式法则（综述）
% license CCBYSA3
% type Wiki

本文根据 CC-BY-SA 协议转载翻译自维基百科\href{https://en.wikipedia.org/wiki/Chain_rule}{相关文章}

在微积分中，链式法则是一条公式，它用函数$f$和$g$的导数来表达它们复合函数的导数。更准确地说，若$h = f \circ g$是函数，满足$h(x) = f(g(x))$对任意 $x$ 成立，则在拉格朗日记号中，链式法则为：
$$
h'(x) = f'(g(x)) \, g'(x).~
$$
或者等价地：
$$
h' = (f \circ g)' = (f' \circ g)\cdot g'.~
$$
链式法则也可以用莱布尼茨记号表示。若变量 $z$ 依赖于变量 $y$，而 $y$ 又依赖于变量 $x$（即 $y$ 和 $z$ 都是依赖变量），那么 $z$ 也通过中间变量 $y$ 依赖于 $x$。在这种情况下，链式法则表示为：
$$
\frac{dz}{dx} = \frac{dz}{dy} \cdot \frac{dy}{dx},~
$$
以及：
$$
\left.\frac{dz}{dx}\right|_{x} = \left.\frac{dz}{dy}\right|_{y(x)} \cdot \left.\frac{dy}{dx}\right|_{x},~
$$
其中竖线记号用来指出导数应当在哪些点进行计算。

在积分中，与链式法则对应的规则是代换积分法。
\subsection{直观解释}
直观上，链式法则表明：如果已知 $z$ 相对于 $y$ 的瞬时变化率，以及 $y$ 相对于 $x$ 的瞬时变化率，那么就可以通过两者的乘积来计算 $z$ 相对于 $x$ 的瞬时变化率。

正如 George F. Simmons 所说：“如果一辆汽车的速度是自行车的两倍，而自行车的速度是步行者的四倍，那么汽车的速度就是步行者的 $2 \times 4 = 8$ 倍。”【1】

这个例子与链式法则的关系如下：设 $z, y, x$ 分别表示汽车、自行车和步行者的位置（变量）。汽车与自行车相对位置的变化率为：$\frac{dz}{dy} = 2$.同样，自行车与步行者相对位置的变化率为：$\frac{dy}{dx} = 4$.因此，汽车与步行者相对位置的变化率为：
$$
\frac{dz}{dx} = \frac{dz}{dy}\cdot \frac{dy}{dx} = 2 \cdot 4 = 8.~
$$
位置变化率就是速度之比，而速度就是位置对时间的导数。即：
$$
\frac{dz}{dx} = \frac{\tfrac{dz}{dt}}{\tfrac{dx}{dt}},~
$$
或者等价地：
$$
\frac{dz}{dt} = \frac{dz}{dx}\cdot \frac{dx}{dt},~
$$
这同样是链式法则的一个应用。
\subsection{历史}
链式法则最早似乎是由戈特弗里德·威廉·莱布尼茨使用的。他利用它来计算$\sqrt{a+bz+cz^{2}}$的导数，把它看作平方根函数与函数$a+bz+cz^{2}$的复合。他在 1676 年的一篇论文中首次提及这一点（其中计算中有一个符号错误）【2】。链式法则的常见记号归功于莱布尼茨【3】。纪尧姆·德·洛必达在其《无穷小分析》中隐含地使用了链式法则。尽管莱昂哈德·欧拉的分析著作写于莱布尼茨发现一百多年之后，但其中并未出现链式法则【需要引证】。一般认为，链式法则的第一个“现代”版本出现在拉格朗日 1797 年的《解析函数论》中；它也出现在柯西 1823 年的《在皇家综合理工学院所作无穷小计算讲义摘要》中【3】。
\subsection{表述}
链式法则最简单的形式适用于实变量的实值函数。它表明：如果 $g$ 在点 $c$ 处可导（即导数 $g'(c)$ 存在），并且 $f$ 在 $g(c)$ 处可导，那么复合函数$f \circ g$在 $c$ 处也是可导的，其导数为【4】：
$$
(f \circ g)'(c) = f'(g(c)) \cdot g'(c).~
$$
该法则有时可简写为：
$$
(f \circ g)' = (f' \circ g) \cdot g'.~
$$
如果 $y = f(u)$，而 $u = g(x)$，那么在莱布尼茨记号中，这种简写形式可写为：
$$
\frac{dy}{dx} = \frac{dy}{du} \cdot \frac{du}{dx}.~
$$
导数计算所对应的点也可以明确写出：
$$
\left.\frac{dy}{dx}\right|_{x=c} = \left.\frac{dy}{du}\right|_{u=g(c)} \cdot \left.\frac{du}{dx}\right|_{x=c}.~
$$
进一步推广，若给定 $n$ 个函数$f_1, f_2, \ldots, f_n$，复合函数为：$f_1 \circ (f_2 \circ \cdots (f_{n-1} \circ f_n))$，若每个函数$f_i$在其输入点上可导，则通过多次应用链式法则，该复合函数也可导，其导数在莱布尼茨记号中为：
$$
\frac{df_1}{dx} = \frac{df_1}{df_2} \cdot \frac{df_2}{df_3} \cdots \frac{df_n}{dx}.~
$$
\subsection{应用}
\subsubsection{两个以上函数的复合}
\begin{figure}[ht]
\centering
\includegraphics[width=6cm]{./figures/d99941d3551df3e8.png}
\caption{} \label{fig_LSFZ_1}
\end{figure}
链式法则可以应用于两个以上函数的复合。要求解多个函数复合的导数，可以注意到：函数$f, g, h$（按此顺序）的复合，其实就是$f$与$g \circ h$的复合。链式法则指出，要计算$f \circ g \circ h$的导数，只需要计算$f$的导数以及$g \circ h$的导数。$f$的导数可以直接计算，而$g \circ h$的导数可以通过再次应用链式法则来计算。

为了更具体，考虑函数：
$$
y = e^{\sin(x^2)}.~
$$
它可以分解为三个函数的复合：
$$
y = f(u) = e^u, \quad u = g(v) = \sin v, \quad v = h(x) = x^2.~
$$
因此：$y = f(g(h(x)))$.

它们的导数分别为：
$$
\frac{dy}{du} = f'(u) = e^u, \quad
\frac{du}{dv} = g'(v) = \cos v, \quad
\frac{dv}{dx} = h'(x) = 2x.~
$$
链式法则表明，在点 $x=a$ 处复合函数的导数为：

$$
\begin{aligned}
(f \circ g \circ h)'(a) 
&= f'((g \circ h)(a)) \cdot (g \circ h)'(a)\\
&= f'((g \circ h)(a)) \cdot g'(h(a)) \cdot h'(a)\\
&= (f' \circ g \circ h)(a) \cdot (g' \circ h)(a) \cdot h'(a)\\
\end{aligned}~
$$
在莱布尼茨记号中，该结果为：
$$
\frac{dy}{dx} = \left.\frac{dy}{du}\right|_{u=g(h(a))} \cdot \left.\frac{du}{dv}\right|_{v=h(a)} \cdot \left.\frac{dv}{dx}\right|_{x=a},~
$$
或者简写为：
$$
\frac{dy}{dx} = \frac{dy}{du} \cdot \frac{du}{dv} \cdot \frac{dv}{dx}.~
$$
因此，导函数为：
$$
\frac{dy}{dx} = e^{\sin(x^2)} \cdot \cos(x^2) \cdot 2x.~
$$
另一种计算该导数的方法是将复合函数 $f \circ g \circ h$ 看作 $(f \circ g)$ 与 $h$ 的复合。以这种方式应用链式法则得到：
$$
\begin{aligned}
(f \circ g \circ h)'(a) &= (f \circ g)'(h(a)) \cdot h'(a)\\
&= f'(g(h(a))) \cdot g'(h(a)) \cdot h'(a)\\
\end{aligned}~
$$
这与上面计算出的结果相同。这是可以预期的，因为：$(f \circ g) \circ h = f \circ (g \circ h)$.

